% Gasflash Equations Report
% Generated on 2026-05-07
% This report summarises the governing equations implemented in the Gasflash CFD solver.

\documentclass[11pt,a4paper]{article}
\usepackage{amsmath,amssymb}
\usepackage{physics}
\usepackage{booktabs}
\usepackage{graphicx}
\usepackage{hyperref}
\usepackage{geometry}
\geometry{margin=1in}
\title{Gasflash – Governing Equations Report}
\author{Gasflash Development Team}
\date{\today}

\begin{document}
\maketitle

\section{Overview}
The Gasflash application solves 1‑D compressible flow problems using a finite‑volume Roe scheme.  The solver can handle a sequence of components (nozzles, ducts, mass‑addition sources, etc.) and includes the following physical models:

\begin{itemize}
  \item Isentropic flow relations
  \item Fanno flow (friction)
  \item Rayleigh flow (heat addition)
  \item Normal shock relations
  \item Mass‑addition (solid grain) with ``only\_mass\_addition`` mode
  \item Real‑gas equation of state (optional)
\end{itemize}

All models are expressed in terms of the primitive variables \(P, T, \rho, M\) and are evaluated at cell interfaces.

\section{Isentropic Relations}
For an ideal gas with ratio of specific heats \(\gamma\):
\begin{align}
  \frac{P}{P_0} &= \left(1 + \frac{\gamma-1}{2}M^2\right)^{-\frac{\gamma}{\gamma-1}}
  \\
  \frac{T}{T_0} &= \left(1 + \frac{\gamma-1}{2}M^2\right)^{-1}
  \\
  \frac{\rho}{\rho_0} &= \left(1 + \frac{\gamma-1}{2}M^2\right)^{-\frac{1}{\gamma-1}}
\end{align}
The critical (choked) Mach number \(M=1\) gives the well‑known pressure ratio:
\begin{equation}
  \left(\frac{P_*}{P_0}\right) = \left(\frac{2}{\gamma+1}\right)^{\frac{\gamma}{\gamma-1}}
\end{equation}
This expression is used in \texttt{compute\_summary} to estimate the choking threshold.

\section{Fanno Flow (Friction)}
Fanno flow accounts for a constant‑area duct with wall friction coefficient \(f\).  The governing relation between Mach number and nondimensional length \(\theta = \frac{4fL}{D}\) is:
\begin{equation}
  \frac{1-M^2}{\gamma M^2} + \frac{\gamma+1}{2\gamma}\ln\!\left[\frac{(\gamma+1)M^2}{2+(\gamma-1)M^2}\right] = \theta + C
\end{equation}
where \(C\) is determined from the inlet conditions.  The solver iteratively solves this equation for \(M\) at each cell.

\section{Rayleigh Flow (Heat Addition)}
For a constant‑area duct with heat addition \(q\) per unit mass, the Rayleigh relations are:
\begin{align}
  \frac{T}{T_0} &= \frac{(\gamma+1)M^2}{(\gamma-1)M^2+2} \cdot \frac{2+(\gamma-1)M^2}{(\gamma+1)M^2}
  \\
  \frac{P}{P_0} &= \left[\frac{(\gamma+1)M^2}{(\gamma-1)M^2+2}\right]^{\frac{\gamma}{\gamma-1}} \cdot \left[\frac{2+(\gamma-1)M^2}{(\gamma+1)M^2}\right]^{\frac{1}{\gamma-1}}
\end{align}
The solver uses these formulas when the component type is ``rayleigh`` and the parameter ``q`` is supplied.

\section{Normal Shock Relations}
Across a normal shock the downstream Mach number \(M_2\) is:
\begin{equation}
  M_2^2 = \frac{1 + \frac{\gamma-1}{2}M_1^2}{\gamma M_1^2 - \frac{\gamma-1}{2}}
\end{equation}
Pressure and temperature jumps are:
\begin{align}
  \frac{P_2}{P_1} &= 1 + \frac{2\gamma}{\gamma+1}(M_1^2-1)
  \\
  \frac{T_2}{T_1} &= \frac{[2\gamma M_1^2-(\gamma-1)]\,[ (\gamma-1)M_1^2+2]}{(\gamma+1)^2M_1^2}
\end{align}
These are called by \texttt{evaluate\_pipeline} when a shock is detected.

\section{Mass‑Addition (Solid Grain)}
The solid‑grain model follows Vieille’s law for burn rate:
\begin{equation}
  r_b = a\,P^n
\end{equation}
with \(a\) and \(n\) taken from the component parameters.  The code now converts pressure from Pa to MPa before applying the law (see \texttt{influence\_solver.py}).  When ``only_mass_addition=True`` the solver skips the aerodynamic coupling and only adds mass to the flow.

\section{Real‑Gas EOS (optional)}
When ``request.is_real`` is true, the ``RealGasProperties`` class uses the Peng‑Robinson EoS:
\begin{equation}
  P = \frac{RT}{v-b} - \frac{a}{v(v+b)+b(v-b)}
\end{equation}
where \(a\) and \(b\) are temperature‑dependent parameters.

\section{Numerical Scheme}
The core solver is a first‑order Roe approximate Riemann solver.  For each interface the flux \(\mathbf{F}\) is computed as:
\begin{equation}
  \mathbf{F}_{i+1/2}=\frac{1}{2}\big[\mathbf{F}(\mathbf{U}_L)+\mathbf{F}(\mathbf{U}_R)\big] - \frac{1}{2}\,\mathbf{R}\,|\Lambda|\,\mathbf{R}^{-1}(\mathbf{U}_R-\mathbf{U}_L)
\end{equation}
with \(\mathbf{U}\) the conserved vector, \(\mathbf{R}\) the eigen‑vector matrix and \(\Lambda\) the diagonal matrix of eigenvalues.  Time integration uses an explicit Euler step with CFL control.

\section{Summary of Changes}
The recent code cleanup removed the legacy analytical (iterative) solver and unified the pipeline under the CFD (hybrid) solver.  All equations above are now evaluated exclusively by the CFD path.

\end{document}
